\begin{proof}
The proof is very similar to that of Thm.~\ref{Thm: Main, not weighted}. After Lemma~\ref{Lemma: weigted enumeration only multiples of |g| contribute}, we need only consider in \eqref{EQ: weighted enum via local Coxeter numbers} those characters $\chi$ for which all $c_{\chi_{,\mathcal{C}}}$ are multiples of $|g|$. This allows us to write the exponential function as $$\op{FAC}_{W,g}(\bm w,z)=\dfrac{e^{z\cdot\op{wt}(\mathcal{R})}}{|W|}\cdot\tilde{\Phi}(\bm X),$$ for a polynomial $\tilde{\Phi}$ in variables $\bm X\coloneqq (X_{\mathcal{C}})_{\mathcal{C}\in\mathcal{A}/W}$, by setting $X_{\mathcal{C}}=\big( e^{-zw_{\mathcal{C}}^{\phantom{l_r}}}\big)^{|g|}$. By Corol.~\ref{Corol: local cox num are integers andbounds} the polynomial $\tilde{\Phi}(\bm X)$ has degree $(e_{\mathcal{C}}\omega_{\mathcal{C}})/|g|$ in each of its variables $X_{\mathcal{C}}$, and it has constant term $1$ since all $c_{\chi_{,\mathcal{C}}}$ can be {\it simultaneously} $0$ only for the trivial representation.

To find the largest power of $(1-X_{\mathcal{C}})$ that divides $\tilde{\Phi}(\bm X)$, we view $\tilde{\Phi}$ as a polynomial in the single variable $X_{\mathcal{C}}$ and treat the other $X_{\mathcal{C}'}$'s as complex scalars. This is equivalent to assigning arbitrary values to all variables $w_{\mathcal{C}'}\neq w_{\mathcal{C}}$ of the weight function in Defn.~\ref{Defn: weight function}. If we further fix $z=1$, the enumerative interpretation of $(e^{\op{wt}(\mathcal{R})}/|W|)\cdot\tilde{\Phi}(X_{\mathcal{C}})$ is then that it counts weighted reflection factorizations of $g$ keeping track only of the number of reflections that fix a hyperplane in $\mathcal{C}$.

Now, as in Thm.~\ref{Thm: Main, not weighted} consider the root factorization of $\tilde{\Phi}(X_{\mathcal{C}})$: $$\tilde{\Phi}(X_{\mathcal{C}})=a(\alpha_1-X_\mathcal{C})(\alpha_2-X_{\mathcal{C}})\cdots (\alpha_r-X_{\mathcal{C}}),$$ with $r=(e_{\mathcal{C}}\omega_{\mathcal{C}})/|g|$. We see again that by plugging back $X_{\mathcal{C}}=e^{-w_{\mathcal{C}}|g|}$ each root contributes a factor of either $(\alpha_i-1)$ or $w_{\mathcal{C}}|g|$ to the leading term of the generating function. Since by Defn.~\ref{Defn: special reflection lengths n_C} this must be a scalar multiple of $w_{\mathcal{C}}^{n_C(g)}$, we have that $(1-X_{\mathcal{C}})^{n_{\mathcal{C}}(g)}$ divides $\tilde{\Phi}(X_{\mathcal{C}})$ and is the largest power that does so (this furthermore proves the first inequality). Since this is true for a dense set of the complex values $X_{\mathcal{C}'}$, we in fact have that $(1-X_{\mathcal{C}})^{n_{\mathcal{C}}(g)}$ is a maximal factor of $\tilde{\Phi}(\bm X)$.

The only thing left to show is the second inequality for the $n_{\mathcal{C}}(g)$'s. To see this, we now identify all weights $w_{\mathcal{C}'},\ \mathcal{C}'\neq\mathcal{C}$ to a single weight $w$, set again $z=1$, and treat $\tilde{\Phi}$ as a polynomial in two variables $X=e^{-w|g|}$ and $X_{\mathcal{C}}=e^{-w^{\phantom{l}}_{\mathcal{C}}|g|}$. The general argument about $\tilde{\Phi}(\bm X)$ implies that we can consider the polynomial $\Phi'(X,X_{\mathcal{C}})$ defined by $$\Phi'(X,X_{\mathcal{C}})\coloneqq \dfrac{\tilde{\Phi}(X,X_{\mathcal{C}})}{(1-X_{\mathcal{C}})^{n_{\mathcal{C}}(g)}}.$$
Now, the generating function $$\dfrac{e^{\op{wt}(\mathcal{R})}}{|W|}\cdot\Phi'(X,X_{\mathcal{C}})\cdot(1-X_{\mathcal{C}})^{n_{\mathcal{C}}(g)}$$ counts 
reflection factorizations of $g$ weighing reflections in $\mathcal{C}^{\op{ref}}$ by $w_{\mathcal{C}}$ and the rest by $w$. 
We want to enumerate factorizations that have exactly the minimal number $n_{\mathcal{C}}(g)$ of reflections of type $\mathcal{C}$. 
Since the term $(1-X_{\mathcal{C}})^{n_{\mathcal{C}}(g)}$ always contributes a factor of $(w^{\phantom{l}}_{\mathcal{C}}|g|)^{n_{\mathcal{C}}(g)}$ to the Taylor expansion, the answer to the previous question would be given by $$\dfrac{|g|^{n_{\mathcal{C}}(g)}}{|W|}\cdot e^{\op{wt}(\mathcal{R})}\Big|_{w_{\mathcal{C}}=0}\cdot\Phi'(X,X_{\mathcal{C}})\Bigg|_{\substack{X_{\mathcal{C}}=1\quad\ \\ X=e^{-w|g|}} }.$$
The leading term of this exponential generating function should clearly be a multiple of $w^{f_{\mathcal{C}}(g)-n_{\mathcal{C}}(g)}$, where $f_{\mathcal{C}}(g)$ is the smallest length of a reflection factorization of $g$ with exactly $n_{\mathcal{C}}(g)$-many reflections of type $\mathcal{C}$. As in the previous argument, this implies that $\Phi'(X,1)$ is a multiple of $(1-X)^{f_{\mathcal{C}}(g)-n_{\mathcal{C}}(g)}$, but since by construction its degree is equal to $\sum_{\mathcal{C}'\neq\mathcal{C}}(e_{\mathcal{C}'}\omega_{\mathcal{C}'})/|g|$, which is in turn equal to $\left(|\mathcal{R}|+|\mathcal{A}|-e_{\mathcal{C}}\omega_{\mathcal{C}}\right)/|g|$, we must have 
\[
f_{\mathcal{C}}(g)-n_{\mathcal{C}}(g)\leq \dfrac{|\mathcal{R}|+|\mathcal{A}|-e_{\mathcal{C}}\omega_{\mathcal{C}}}{|g|},
\]
 which completes the proof, since $f_{\mathcal{C}}(g)\geq l_R(g)$.
\end{proof}
